\documentclass[11pt]{article}
\usepackage[T1]{fontenc}
\usepackage[utf8]{inputenc}
\usepackage[margin=1in]{geometry}
\usepackage{amsmath,amssymb}

\setlength{\parindent}{0pt}
\setlength{\parskip}{0.6em}

\begin{document}

\title{Technical Review of ``MSAGAT-Net: Multi-Scale Temporal Adaptive Graph Attention for Efficient Spatiotemporal Epidemic Forecasting''}
\author{}
\date{2026-04-09}
\maketitle

\section*{Overall assessment}

This is a technically ambitious manuscript with several genuinely interesting design ideas, especially the optional additive adjacency prior, the graph-size-adaptive spatial aggregation, and the cross-dataset empirical scope. I explicitly checked symbol definitions, dimensional consistency, sign/branch-convention issues, figure/table-to-text alignment, and implementation-facing formulas. I did not find inverse-trigonometric or sign-convention errors, but I did find one definite mathematical problem in the attention regulariser, one definite arithmetic inconsistency in the adaptive hop-depth rule, one major loss-definition inconsistency, and several places where the manuscript states stronger claims than the presented evidence supports.

In its current form, the paper reads as promising but not yet technically settled. The highest-priority revisions are mathematical/implementation-facing rather than stylistic.

\section*{Highest-priority fixes}

\begin{enumerate}
\item Replace or reformulate the post-softmax attention sparsity penalty.
\item Make the training objective consistent across the attention-regularisation subsection, the total-loss equation, Table~2, and Algorithm~1.
\item Reconcile the adaptive hop-depth formula with the reported dataset-specific hop counts; the current text is arithmetically inconsistent for the 10-node US-Region dataset.
\item State the temporal train/validation/test split protocol explicitly enough to rule out forecasting leakage.
\end{enumerate}

\section*{Technical findings}

\begin{enumerate}

\item \textbf{Definite error: the attention regularisation term is ineffective as written.}

\textbf{Location:} subsection ``Attention Regularisation'', especially
\[
\mathcal{L}_{\text{attn}} = \lambda \lVert \mathbf{A}_h \rVert_1.
\]

\textbf{Problem:} the manuscript defines $\mathbf{A}_h$ as the row-wise softmax output used for value aggregation. If $\lVert \cdot \rVert_1$ is the elementwise $L_1$ norm (which is what the accompanying prose strongly suggests), then each row of $\mathbf{A}_h$ is nonnegative and sums to $1$, so the penalty is constant after softmax and cannot promote row sparsity.

\textbf{Why it matters:} the paper currently presents this term as a mechanism that ``directly encourages each node to attend to a sparse subset of other nodes''. As written, that claim is not mathematically correct, and the regulariser does not do the advertised job.

\textbf{Suggested fix:} either (i) redefine the sparsity term using a non-constant surrogate such as row-wise entropy, sparsemax/entmax, or a pre-softmax penalty on the structural-bias parameters, or (ii) if the induced matrix $1$-norm is intended, state that explicitly and rewrite the explanation because the current interpretation is about row sparsity rather than maximum column mass.

\item \textbf{Definite inconsistency: the attention-loss weighting is defined three different ways.}

\textbf{Location:} the ``Attention Regularisation'' subsection, the total-loss equation in the training section, Table~2, and Algorithm~1.

\textbf{Problem:} the paper first defines $\mathcal{L}_{\text{attn}} = \lambda \lVert \mathbf{A}_h \rVert_1$, then later writes
\[
\mathcal{L}_{\text{total}} = \mathcal{L}_{\text{forecast}} + \lambda_{\text{attn}}\mathcal{L}_{\text{attn}},
\]
while Algorithm~1 uses $\mathcal{L}_{\text{MSE}}(\hat{\mathbf{Y}}, \mathbf{y}) + \mathcal{L}_{\text{attn}}$ with no outer factor. As written, the regulariser is weighted once, weighted twice, or notated under two different symbols depending on which part of the manuscript the reader follows.

\textbf{Why it matters:} this is a reproducibility problem. An implementer cannot recover the intended optimisation objective unambiguously.

\textbf{Suggested fix:} choose one convention and propagate it everywhere. For example, define $\mathcal{L}_{\text{attn}}$ without its coefficient, then use $\mathcal{L}_{\text{total}} = \mathcal{L}_{\text{forecast}} + \lambda_{\text{attn}}\mathcal{L}_{\text{attn}}$ consistently in the text, table, and pseudocode.

\item \textbf{Definite arithmetic inconsistency: the reported hop counts do not follow the stated rule.}

\textbf{Location:} subsection ``Adaptive Hop Depth and Locality-Biased Fusion''.

\textbf{Problem:} the manuscript states
\[
S = \min(S_{\max}, \max(2, \lfloor N/5 \rfloor)),
\]
and then says this yields $S=2$ for Australia-COVID and NHS-Timeseries, and $S=4$ for all other datasets. However, Table~1 reports that the US-Region dataset has $N=10$, so the stated formula gives
\[
S = \min(4, \max(2, \lfloor 10/5 \rfloor)) = 2,
\]
not $4$.

\textbf{Why it matters:} this is a direct mismatch between the documented algorithm and the dataset-level description.

\textbf{Suggested fix:} either correct the formula, or correct the sentence describing the realised hop counts. A small table listing the actual $S$ used for each dataset would remove the ambiguity entirely.

\item \textbf{Likely issue / implementation ambiguity: the split description is not sufficient to rule out temporal leakage.}

\textbf{Location:} subsection ``Training and Optimisation Strategy''.

\textbf{Problem:} the manuscript states a $60\%{:}20\%{:}20\%$ train/validation/test split and mentions a fixed random seed, but it does not say whether the split is chronological, whether validation/test windows are strictly later in time than training windows, or whether sliding windows are prevented from crossing split boundaries.

\textbf{Why it matters:} for forecasting tasks, random or improperly windowed splits can leak future information into training and materially inflate test performance.

\textbf{Suggested fix:} state explicitly that the split is chronological (if that is what was done), give the exact temporal cut points, and clarify how sliding-window samples are generated near split boundaries.

\item \textbf{Likely issue: EAGAM mixes per-head and stacked-head notation.}

\textbf{Location:} subsection ``Multi-Head Attention Mechanism''.

\textbf{Problem:} the text says ``within each head we compute''
\[
\mathbf{S}_h = \frac{\mathbf{Q}_h \mathbf{K}_h^T}{\sqrt{d_{\text{head}}}},
\]
but then assigns $\mathbf{S}_h \in \mathbb{R}^{h \times N \times N}$. The same overload appears for $\mathbf{A}_h$ and $\mathbf{O}_h$.

\textbf{Why it matters:} this conflates a specific head index with a stacked tensor over all heads, which makes the implementation-facing math harder to follow and increases the chance of reader error.

\textbf{Suggested fix:} use a separate symbol for the stacked tensor, or index a single head explicitly, e.g. $\mathbf{S}^{(m)} \in \mathbb{R}^{N \times N}$ for head $m$ and $\mathbf{S} \in \mathbb{R}^{h \times N \times N}$ for the stacked object.

\item \textbf{Unsupported claim: the self-regulating prior is described more strongly than the current math proves.}

\textbf{Location:} abstract, subsection ``Self-Regulating Additive Adjacency Prior'', Figure~3 caption, and Conclusion.

\textbf{Problem:} the manuscript repeatedly says that the adjacency prior's influence ``adapts automatically with graph density'' and can become effectively invisible on dense graphs by softmax shift-invariance. Shift-invariance cancels exact row-wise constant offsets; it does not imply that every dense row-normalised adjacency prior is harmless. A dense but non-uniform prior can still perturb the attention distribution materially.

\textbf{Why it matters:} the current wording turns an empirical tendency into a general guarantee.

\textbf{Suggested fix:} qualify the claim to something like: when the row-normalised prior is close to uniform, its additive effect is attenuated by softmax shift-invariance. If a stronger claim is intended, add either a formal bound or an empirical diagnostic based on row-wise variance of $\tilde{\mathbf{A}}$.

\item \textbf{Unsupported explanatory claims: several result interpretations need either evidence or softer wording.}

\textbf{Location:} the paragraphs interpreting learned attention and the Australia-COVID results.

\textbf{Problem:} the manuscript attributes learned attention structure to ``commuter flows, air travel routes, and socioeconomic links'', and attributes the Australia-COVID behaviour to border closures ``effectively decoupling spatial transmission''. These are plausible interpretations, but they are not demonstrated by the experiments currently shown.

\textbf{Why it matters:} readers may take these explanations as established findings rather than hypotheses.

\textbf{Suggested fix:} either support these claims with external mobility/policy evidence or rephrase them as plausible interpretations rather than conclusions.

\end{enumerate}

\section*{Meaningful editorial findings}

\begin{enumerate}

\item \textbf{Location:} results overview and influenza performance discussion.

\textbf{Problem:} the weekly influenza datasets are repeatedly described as ``3, 5, 10, and 15 days ahead'', even though Table~1 marks Japan, US-Region, and US-State as \emph{weekly} datasets.

\textbf{Why it matters:} the forecast-horizon unit is wrong.

\textbf{Suggested fix:} change these references to ``steps ahead'' or ``weeks ahead'' for the weekly datasets, and reserve ``days'' for the daily COVID/NHS datasets.

\item \textbf{Location:} dataset names across Table~1, the results tables, MSSFM discussion, and the conclusion.

\textbf{Problem:} the LTLA dataset appears under three names (LTLA-COVID, LTLA-TimeSeries, LTLA-Timeseries) and the NHS dataset also appears under three names (NHS-ICUBeds, NHS-TimeSeries, NHS-Timeseries). Similar singular/plural drift appears for Japan-Prefecture(s), US-Region(s), and US-State(s).

\textbf{Why it matters:} the naming drift makes it look as though multiple datasets are being discussed.

\textbf{Suggested fix:} choose one canonical name per dataset and use it everywhere, including tables, figures, captions, and the conclusion.

\item \textbf{Location:} influenza dataset subsection.

\textbf{Problem:} the official name ``Centers for Disease Control and Prevention'' is written as ``Centres''.

\textbf{Why it matters:} proper names should preserve official spelling even in UK-English prose.

\textbf{Suggested fix:} change this to ``Centers for Disease Control and Prevention (CDC)''.

\item \textbf{Location:} NHS-Timeseries results paragraph.

\textbf{Problem:} the sentence ``EpiGNN leads at 14-day (0.869 vs 0.869, effectively tied)'' contradicts the table as written. The rounded PCC values shown in the table are tied, not a lead.

\textbf{Why it matters:} the prose overstates a comparison that the reported numbers do not support.

\textbf{Suggested fix:} write ``tied to three decimals'' or report more precision if there is a genuine lead.

\item \textbf{Location:} spatial graph construction / preprocessing description.

\textbf{Problem:} the 7-day rolling-mean smoother is described without saying explicitly whether it applies only to the daily datasets or also to the weekly influenza series.

\textbf{Why it matters:} this is a preprocessing ambiguity, and applying a 7-day smoother to weekly data would be nonsensical.

\textbf{Suggested fix:} state the preprocessing pipeline separately for weekly and daily datasets.

\item \textbf{Location:} bibliography entry \texttt{lijingwangCausalGNNCausalBasedGraph2022} in \texttt{ref.bib}.

\textbf{Problem:} the author list contains duplicated names and the venue field reads ``Proceedings of the ... AAAI Conference on Artificial Intelligence''.

\textbf{Why it matters:} this is an obvious bibliography-import artifact that reduces professionalism and may distort the rendered reference.

\textbf{Suggested fix:} deduplicate the author list and replace the venue field with the correct conference/journal title string.

\end{enumerate}

\section*{Proofing sweep}

\begin{itemize}
\item I ran the mandated micro-proofing scan via \texttt{scripts/proofing\_scan.py}. Most automatic hits were false positives caused by mathematical ellipses (for example $x_1, x_2, \ldots, x_T$) or bibliography access-date notes, so I would not elevate those to manuscript errors.
\item The prose around the aggregated-prediction figure refers to ``confidence bands'', but the caption describes ``$\pm 1$ standard deviation bands''. If these are empirical dispersion bands rather than predictive uncertainty intervals, ``confidence'' is the wrong term. Use ``standard-deviation bands'' or ``dispersion bands'' consistently.
\item The manuscript mixes \texttt{TimeSeries} and \texttt{Timeseries}. Standardising one spelling throughout would remove a visible copyediting distraction.
\end{itemize}

\section*{Items needing external verification}

\begin{itemize}
\item If the manuscript wishes to retain the interpretation that learned attention captures commuter-flow / air-travel / socioeconomic pathways, that claim needs external corroboration rather than qualitative intuition alone.
\item The explanation that Australia's border-closure regime suppressed spatial coupling is plausible, but the current manuscript does not test it directly; either support it externally or soften it.
\item The Data Availability statement claims that all datasets and code are publicly available in the repository. That should be verified against the public repository state before submission.
\end{itemize}

\end{document}